\chapter{Finite-Dimensional Operator Foundations}
\index{operator!finite-dimensional foundations|textbf}

\label{ch:finite-dimensional-foundations}

\chapterstatus{proof contracts and proof strategies are presented with explicit status; no proposed physics proof package is represented as complete or reviewed.}


\section{Finite-dimensional retained frames}
\index{frame!retained}


Let $V\in\mathbb C^{D\times M}$ have orthonormal columns,
\begin{equation}
  V^{\dagger}V=I_M,
\end{equation}
and let $G\in\mathbb C^{M\times M}$ be unitary. The columns of $V$ represent an
ordered orthonormal frame for an $M$-dimensional retained subspace of an ambient
$D$-dimensional space.

\begin{proposition}[\texttt{PRF-05.01}: projector invariance]
\label{prop:projector-invariance}
For $P=VV^{\dagger}$ and $V'=VG$, the rotated frame induces the same ambient
orthogonal projector:
\begin{equation}
  V'V'^{\dagger}=P.
\end{equation}
\end{proposition}

\begin{proof}
Using $(VG)^{\dagger}=G^{\dagger}V^{\dagger}$ and
$GG^{\dagger}=I_M$,
\begin{equation}
  (VG)(VG)^{\dagger}
  =VGG^{\dagger}V^{\dagger}
  =VV^{\dagger}.
\end{equation}
The condition $V^{\dagger}V=I_M$ is needed to interpret $VV^{\dagger}$ as the
orthogonal projector, although the displayed equality uses only the unitarity
of $G$.
\end{proof}

This is the sole frozen contract currently encoded and checked in Lean. The
pinned trial uses Lean~4 and mathlib \texttt{v4.33.0}; the project source contains
no \texttt{sorry} or \texttt{admit}. The check establishes acceptance of that
encoding, not physical admissibility of the frame, smoothness in $\mathbf k$,
or conformance with Isabelle or Rocq.

\citationtodo{medium priority}{Check and add system references for Lean 4,
mathlib, Isabelle/HOL, and Rocq. Provisional keys:
\texttt{demouraUllrich2021}, \texttt{mathlib2020},
\texttt{nipkowPaulsonWenzel2002}, and \texttt{bertotCasteran2004}. The claim that
this repository contains a checked encoding and no \path{sorry} or
\path{admit} still requires durable internal evidence rather than scholarly
citations.}

\begin{proposition}[\texttt{PRF-05.02}: covariance of compression]
\label{prop:compression-covariance}
For an ambient operator $A\in\mathbb C^{D\times D}$, define
\begin{equation}
  H=V^{\dagger}AV,
  \qquad
  H'=(VG)^{\dagger}A(VG).
\end{equation}
Then
\begin{equation}
  H'=G^{\dagger}HG.
\end{equation}
\end{proposition}

\begin{proof}
Associativity and the adjoint product rule give
\begin{equation}
  (VG)^{\dagger}A(VG)
  =G^{\dagger}(V^{\dagger}AV)G.
\end{equation}
No Hermiticity assumption on $A$ is needed for this coordinate identity.
\end{proof}

The contract does not assert that $A$ preserves the retained subspace or that
the compressed spectrum equals the full ambient spectrum.

\section{Identification pullback and aligned subtraction}
\index{pullback!identification}
\index{aligned subtraction}


Let $H_b\in\mathbb C^{M_b\times M_b}$ and
$H_d\in\mathbb C^{M_d\times M_d}$ represent pristine and doped retained
operators. Let
\begin{equation}
  U:\mathbb C^{M_b}\longrightarrow\mathbb C^{M_d}
\end{equation}
be an identification map. Under independent unitary frame changes $G_b$ and
$G_d$, define
\begin{equation}
  H_b'=G_b^{\dagger}H_bG_b,
  \qquad
  H_d'=G_d^{\dagger}H_dG_d,
  \qquad
  U'=G_d^{\dagger}UG_b.
\end{equation}

\begin{proposition}[\texttt{PRF-05.03}: pullback covariance]
\label{prop:pullback-covariance}
The doped operator pulled into pristine coordinates transforms as
\begin{equation}
  U'^{\dagger}H_d'U'
  =
  G_b^{\dagger}(U^{\dagger}H_dU)G_b.
\end{equation}
\end{proposition}

\begin{proof}
Insert the primed definitions and cancel
$G_dG_d^{\dagger}=I_{M_d}$ on both sides of $H_d$:
\begin{align}
  U'^{\dagger}H_d'U'
  &=G_b^{\dagger}U^{\dagger}G_d
    G_d^{\dagger}H_dG_d
    G_d^{\dagger}UG_b\\
  &=G_b^{\dagger}U^{\dagger}H_dUG_b.
\end{align}
Unitarity of $U$ is not required for this covariance identity.
\end{proof}

Two common-space differences must be distinguished:
\begin{align}
  \Delta H_{d\to b}
  &=U^{\dagger}H_dU-H_b,
  &\Delta H_{d\to b}&:\mathbb C^{M_b}\to\mathbb C^{M_b},\\
  \Delta H_{b\to d}
  &=H_d-UH_bU^{\dagger},
  &\Delta H_{b\to d}&:\mathbb C^{M_d}\to\mathbb C^{M_d}.
\end{align}

\begin{proposition}[\texttt{PRF-05.04}: aligned-difference covariance and equivalence]
\label{prop:aligned-difference-equivalence}
Each difference transforms covariantly in its own common space. If $U$ is a
unitary identification, then
\begin{equation}
  \Delta H_{b\to d}
  =U\Delta H_{d\to b}U^{\dagger},
  \qquad
  \Delta H_{d\to b}
  =U^{\dagger}\Delta H_{b\to d}U.
\end{equation}
\end{proposition}

\begin{proof}
Covariance follows by applying Proposition~\ref{prop:pullback-covariance} to the
pullback difference and its pushforward analogue to the doped-space difference.
For unitary $U$,
\begin{align}
  U\Delta H_{d\to b}U^{\dagger}
  &=U(U^{\dagger}H_dU-H_b)U^{\dagger}\\
  &=H_d-UH_bU^{\dagger}
  =\Delta H_{b\to d}.
\end{align}
The inverse relation follows similarly.
\end{proof}

This algebra does not establish that a physically adequate $U$ exists, that it
is smooth or symmetry compatible, or that the scalar energy zeros of $H_b$ and
$H_d$ agree. Those remain prerequisites before either difference is called an
impurity operator.

\section{Invariant norms and gauge-dependent locality}
\index{norm!unitary invariant}
\index{locality!gauge dependence}


For a unitary $G$, the Frobenius and induced Euclidean operator norms satisfy
\begin{equation}
  \|G^{\dagger}AG\|_{\mathrm F}=\|A\|_{\mathrm F},
  \qquad
  \|G^{\dagger}AG\|_2=\|A\|_2.
\end{equation}
These are the frozen \texttt{PRF-05.05a} and \texttt{PRF-05.05b} targets. They
support gauge-invariant residuals only when the compared paths transform
conformantly, as required by \texttt{PRF-05.06}.

Spatial blocks and shell decompositions require a narrower statement. A general
unitary rotation can mix orbitals associated with different sites or cells and
therefore change which entries are classified as local or exterior. The frozen
\texttt{PRF-05.07} target is a two-point counterexample showing that gauge
transformation and real-space truncation need not commute. The
\texttt{PRF-05.08} target instead restricts the gauge to declared lattice-local
constant rotations under which shell Frobenius diagnostics remain invariant.

Thus a global unitarily invariant norm and a spatially resolved residual have
different admissible gauge groups. A proof about the first cannot be reused for
the second without the additional locality assumptions.

\section{TB-anchored identification}
\index{alignment!tight-binding anchored}


Suppose $\widetilde V_b(\mathbf k)$ and
$\widetilde V_d(\mathbf k)$ are orthonormal frames of equal rank after both have
been related to a declared tight-binding anchor. The candidate fiberwise map is
\begin{equation}
  U_d(\mathbf k)
  =
  \widetilde V_d(\mathbf k)
  \widetilde V_b(\mathbf k)^{\dagger}.
\end{equation}
On the corresponding retained spaces, this formula gives an isometric unitary
identification under the stated orthonormality and equal-rank assumptions.

The substantial \texttt{PRF-10.01} obligations are not the elementary product
identity. They concern existence of suitable anchored frames, rank and
conditioning, global smoothness and periodicity, crystal symmetry, and physical
adequacy. \texttt{PRF-10.02} then owns equivariance of the construction, while
\texttt{PRF-10.03} owns the common-space subtraction and its covariance. All
three remain proposed.
